\section{Diameter directions and the discrete Borsuk partition problem} \label{sec:directions-Borsuk}

We begin with \Cref{thm:diam-dir} which bounds the number of diameter directions of a lattice polytope. Thereafter, we expand our study of lattice diameters to bounded sets $S \subset \Z^d$. Note that if $S = \conv(S) \cap \Z^d$, then $\ldiam(S)=\ldiam(\conv (S))$, where $\conv (S)$ is a lattice polytope, but our results hold for any discrete set of lattice points $S$. In \Cref{Lemma:MaxDegree}, we prove an upper bound on the number of lattice diameter segments that end at a given point in $S$. Using this, we prove \Cref{Thm:Borsuk}, the discrete Borsuk partition problem. A lemma which will be of much use in this section, from \cite{rabi89}, is the following:

\begin{lemma}[Rabinowitz, 1989]\label{lemma:Rabinowitz}
    Let $S \subseteq \Z^d$ be a bounded set and let $m \in \N$ satisfy $|S| > m^d$. Then $\conv (S)$ has at least $m+1$ collinear points. In other words, if $\ldiam(S)<m$, then $|S| \leq  m^d$.
\end{lemma}

% \begin{proof}
%     Suppose, for the sake of contradiction, that $\lvert S \rvert > 2^d$. Each point in $\Z^d$ has a well-defined parity vector in $\{0,1\}^d$, obtained by reducing each coordinate modulo 2. Since there are only $2^d$ such parity vectors, by the pigeonhole principle, there exist two distinct points $x, y \in S$ with the same parity vector. This implies that $x \equiv y \pmod{2}$, so their average $(x+y)/2$ is an integer point. But then the midpoint lies in $\Z^d$ and is distinct from both $x$ and $y$, contradicting the assumption that $\diam_\Z(S) = 1$, which means no lattice point lies strictly between any two points of $S$. Therefore, we must have $\lvert S \rvert \leq 2^d$.
% \end{proof}

\subsection{Diameter directions of a lattice polytope}

Alarcon \cite[Theorem~4.1]{Alarcon_PhD} showed that the number of diameter directions of a lattice polygon $P$ is at most six if $\ldiam(P)=1$ and at most four if $\ldiam(P)>1$; see \Cref{fig:diam-dir} for an example.


\begin{figure}[ht!]
  \centering
  % Figures/3.1_diam-dir.tex
% (Se asume \usepackage{tikz} en el preámbulo)
% Definiciones de colores si no existen (para evitar "Undefined color")
% \makeatletter
% \@ifundefined{color@lilac}{\definecolor{lilac}{RGB}{200,162,200}}{}
% \@ifundefined{color@darklav}{\definecolor{darklav}{RGB}{102,51,153}}{}
% \@ifundefined{color@lightgreen}{\definecolor{lightgreen}{RGB}{102,187,106}}{}
% \@ifundefined{color@stronggreen}{\definecolor{stronggreen}{RGB}{27,94,32}}{}
% \@ifundefined{color@darkyellow}{\definecolor{darkyellow}{RGB}{184,134,11}}{}
% \@ifundefined{color@stoneblue}{\definecolor{stoneblue}{RGB}{100,149,237}}{}
% \makeatother

\begin{minipage}{0.48\textwidth}
  \centering
  \begin{tikzpicture}[thick, scale=0.7]
    % malla
    \foreach \x in {-1,...,3} {
      \foreach \y in {-1,...,3} {
        \node at (\x,\y) [circle, draw=lightgray, fill=lightgray, minimum width=2pt] {};
      }
    }
    % polígono
    \draw[line width=0.3mm, gray, fill=lightgray, opacity=.5] (0,1) -- (1,0) -- (2,2) -- cycle;

    % rectas (recortadas al cuadro [-1,3]x[-1,3])
    \draw[line width=0.4mm, orange] (-1,1) -- (3,1);                 % horizontal
    \draw[line width=0.4mm, lilac!50!darklav] (1,-1) -- (1,3);        % vertical
    \draw[line width=0.4mm, lightgreen] (-1,-1) -- (3,3);             % SW-NE
    \draw[line width=0.4mm, stronggreen] (2,-1) -- (-1,2);            % SE-NW
    \draw[line width=0.4mm, darkyellow] (-1,0.5) -- (3,2.5);          % paralela arriba
    \draw[line width=0.4mm, stoneblue] (0.5,-1) -- (2.5,3);           % paralela abajo

    % puntos negros
    \foreach \p in {(1,0),(0,1),(2,2),(1,1)}{
      \node at \p [circle, fill=black, minimum width=3.5pt] {};
    }
  \end{tikzpicture}
    \hspace{-8em}
    \end{minipage} 
    \hfill
\begin{minipage}{0.48\textwidth}
    \hspace{-8em}
  \centering
  \begin{tikzpicture}[thick, scale=0.7]
    % malla
    \foreach \x in {-1,...,3} {
      \foreach \y in {-1,...,3} {
        \node at (\x,\y) [circle, draw=lightgray, fill=lightgray, minimum width=2.5pt] {};
      }
    }
    % polígono
    \draw[line width=0.3mm, gray, fill=lightgray, opacity=0.5] (0,0) -- (2,0) -- (2,2) -- (0,2) -- cycle;

    % rectas hasta la frontera
    \draw[line width=0.4mm, orange] (-1,0) -- (3,0);
    \draw[line width=0.4mm, orange] (-1,1) -- (3,1);
    \draw[line width=0.4mm, orange] (-1,2) -- (3,2);
    \draw[line width=0.4mm, lilac!50!darklav] (0,-1) -- (0,3);
    \draw[line width=0.4mm, lilac!50!darklav] (1,-1) -- (1,3);
    \draw[line width=0.4mm, lilac!50!darklav] (2,-1) -- (2,3);
    \draw[line width=0.4mm, lightgreen] (-1,-1) -- (3,3);
    \draw[line width=0.4mm, stronggreen] (3,-1) -- (-1,3);

    % puntos negros
    \foreach \p in {(2,0),(0,2),(2,2),(0,0),(1,0),(0,1),(1,1),(2,1),(1,2)}{
      \node at \p [circle, fill=black, minimum width=3.5pt] {};
    }
  \end{tikzpicture}
\end{minipage}

  \caption{Lattice diameter lines colored by their diameter direction.}
  \label{fig:diam-dir}
\end{figure}


We show an analogous statement to the first part, when $\ldiam(P)=1$, in higher dimensions. Our computational experiments suggest that an even smaller bound holds when $\ldiam(P)>1$, but this remains an open problem.

Define the direction $\dir(\bx,\by)$ between two points $\bx,\by \in \R^d$ as $\by-\bx$ or any of its non-zero multiples. For $X \subseteq \R^d$ (usually discrete) define the set of its directions as
\begin{align*}
    \dir(X):= \lb \dir(\bx_1,\bx_2): \bx_1,\bx_2 \in X, \ \bx_1 \neq \bx_2  \rb.
\end{align*}

\begin{lemma}\label{lemma:hyp-dir}
    Let $X,Y \subseteq \R^d$ be sets contained in parallel hyperplanes $H(\ba,\alpha)$ and $ H(\ba,\beta)$, respectively. Suppose that $\dir(X) \cap \dir(Y) = \emptyset$. Then all $\dir(\bx,\by)$ where $\bx \in X$ and $\by \in Y$ are \mbox{pairwise distinct}. 
\end{lemma}

\begin{proof}
    Suppose there exist two distinct pairs $(\bx_1,\by_1), (\bx_2,\by_2) \in X \times Y$ with $\dir(\bx_1,\by_1)=\dir(\bx_2,\by_2)$, that is, $\bx_1-\by_1=\lambda(\bx_2-\by_2)$ for some non-zero $\lambda \in \R$. Then, since $\bx_i \in X \subseteq H(\ba,\alpha)$ and $\by_i \in Y \subseteq H(\ba,\beta)$ we have
    \begin{align*}
        0 \neq \alpha - \beta = \langle a, \bx_1 - \by_1 \rangle = \langle a, \lambda(\bx_2 - \by_2) \rangle 
        = \lambda \langle \ba, \bx_2 - 
        \by_2 \rangle 
        =\lambda (\alpha - \beta)
    \end{align*}
    and thus $\lambda =1$. It follows that $\bx_1-\bx_2 =\by_1-\by_2 \in \dir(X) \cap \dir(Y)$, a contradiction.
\end{proof}

\begin{theorem}\label{thm:diam-dir}
    Let $P$ be a lattice $d$-polytope, with $\ldiam(P)=1$. Then $P$ has at most $\binom{2^d}{2}$ diameter directions and this bound is best possible.
\end{theorem}

\begin{proof}
Let $d \in \N$. For $d=1$ the claim is clear, so assume $d \geq 2$. Since $\ldiam(P)<2$, by \ref{lemma:Rabinowitz} $|P \cap \Z^d| \leq 2^d$. A lattice diameter line is determined by two lattice points in $P$, so the number of possible diameter directions is bounded above by $\binom{2^d}{2}$.

Next, we give a construction of a lattice $d$-polytope attaining this bound. For $t \in \N$ and $x \in \N_{\geq 3}$, let\vspace{-5pt}
    \begin{align*}
        \bv_1=(1,t), \ \bv_2=(x, tx+1), \ \bv_3=(-x-1, -tx-t-1)
    \end{align*}
    and let $T_{t,x} = \conv (\{\bv_1,\bv_2,\bv_3\})$.
    Then $(0,0) = \frac{1}{3}\bv_1+ \frac{1}{3}\bv_2+ \frac{1}{3}\bv_3 \in T_{t,x}$ and since  $T_{t,x}$ has area $\frac{3}{2}$, Pick's Theorem \cite{pick99} implies that $T_{t,x}$ contains no lattice points besides $(0,0)$. Moreover, no three points in $T_{t,x}$ are collinear, thus $\ldiam(T_{t,x})=1$. The slopes of the six diameter lines in $T_{t,x}$ are
    \begin{align*}
    t, \ t + \dfrac{1}{2x + 1}, \ t + \dfrac{1}{x + 2}, \ t + \dfrac{1}{x + 1}, \ t + \dfrac{1}{x}, \text{ and }\ t + \dfrac{1}{x - 1},
    \end{align*}
    which are all distinct and contained in $[t,t+\frac{1}{2}]$.

    %We prove via induction over $k$, for $2 \leq k \leq d$, that there exists a lattice $k$-polytope with $\ldiam(P)=2$ that has $\binom{2^k}{2}$ diameter directions. For $d=k$ this proves the claim.
    \emph{Induction Claim:} For $k \in \{2, \ldots, d\}$ there exist $2^{d-k}$ lattice $k$-polytopes $P_1, \ldots, P_{2^{d-k}}$ such that $|P_i \cap \Z^k| = 2^k$, $P_i$ has $\binom{2^k}{2}$ diameter directions, and $\dir(P_i) \cap \dir(P_j)  = \emptyset$ for all $i,j \in [2^{d-k}], \ i \neq j$.
    
    \emph{Base Case:} For $k=2$ consider $\{ T_{i,3} \st i=1, \ldots, 2^{d-2} \}$. As we showed above $T_{i,3}$ contains $4=2^2$ lattice points, $\ldiam(T_{i,3})=2$ and all lines through two different lattice points have a different diameter direction, thus $T_{i,3}$ has $\binom{4}{2}=6$ diameter directions. Furthermore $\dir(T_{i,3}) \subseteq [i,i+1/2]$, thus the diameter directions of all $T_{i,3}$ for $i =1, \ldots, 2^{d-2}$  are distinct. 

    \emph{Induction Hypothesis:} For $k \in \{2, \ldots, d-1\}$ let $P_1, \ldots, P_{2^{d-k}}$ be lattice $k$-polytopes such that $|P_i \cap \Z^k| = 2^k$, $P_i$ has $\binom{2^k}{2}$ diameter directions, and $\dir(P_i) \cap \dir(P_j)  = \emptyset$ for all $i,j \in [2^{d-k}], \ i \neq j$. In particular, any two pairs of lattice points in $P_i$ define lattice diameter lines \mbox{with distinct slopes.}

    \emph{Induction Step:} Let $P_i$ for $i \in [2^{d-k}]$ be as given in the induction hypothesis. Let $\underline{P_i}:= P_{2i-1} \times \{0\}$ and $\overline{P_{2i}}:= P_{2i} \times \{1\}$ and construct polytopes
    \begin{align*}
        Q_i:= \conv \lp \lb \underline{P_i},\overline{P_{2i}}  \rb \rp\subseteq \R^{k+1}.
    \end{align*}
    Since $\underline{P_i}$ and $\overline{P_{2i}}$ are contained in parallel hyperplanes
    \begin{align*}
        |Q_i \cap \Z^{k+1}| = |\underline{P_i} \cap \Z^{k+1}| +|\overline{P_{2i}} \cap \Z^{k+1}| = 2^{k+1}, 
    \end{align*}
    and $\ldiam(Q_i)=2$. By \Cref{lemma:hyp-dir} all $\dir(\bx,\by)$ for any $\bx \in \underline{P_i}$ and $\by \in \overline{P_{2i}}$ are distinct so any two points in $Q_i$ define different lattice diameter lines with distinct diameter directions, i.e., $Q_i$ has $\binom{2^{k+1}}{2}$ diameter directions. Furthermore, again by \Cref{lemma:hyp-dir}, any two diameter directions in $Q_i$ and $Q_j$ are distinct.
    This proves the induction claim. 
\end{proof}

\begin{example}
    We give an example of our recursive construction using the same idea as in the proof of \mbox{\Cref{thm:diam-dir}}.
    Here $T_0=\conv(\{(0,-3), (1,-4), (5,-5)\})$ and $T_1 = \conv(\{ (3,-3),(4,-2), (2,-7)\})$ each have $6=\binom{2^2}{2}$ distinct diameter directions and thus $P = \conv( T_0 \times \{0\}, T_1 \times \{1\})$ has $\binom{2^3}{2}$ diameter directions. See \Cref{fig:construction-diam-dir}.
\begin{figure}[ht!]
    \centering
    \tikzset{every picture/.style={line width=0.75pt}} %set default line width to 0.75pt        

\begin{tikzpicture}[x=0.9pt, y=0.9pt, yscale=-.7, xscale=.8, transform shape=false]
\begin{scope}
    \clip (0,74) rectangle (350,210);
%uncomment if require: \path (0,299); %set diagram left start at 0, and has height of 299

% % x-axis (behind)
% \draw[color=black] (12,161) -- (355,181);

% % z-axis bottom, dashed inside and top
% \draw[color=black] (190,181) -- (190,208);
% \draw[color=black, dash pattern={on 2.25pt off 2.25pt}] (190,113) -- (190,181);
% \draw[color=black] (190,75) -- (190,113);

% % y-axis
% \draw [color=black] (46,183) -- (331,159);

% x-axis (slope = 1/17), choose x so x-190 is multiple of 17 ⇒ integer y
\draw[color=black] (20,161) -- (360,181);

% y-axis (slope = -1/12), choose x so x-190 is multiple of 12 ⇒ integer y
\draw[color=black] (46,183) -- (346,158);

% z-axis (vertical through the origin)
\draw[color=pink, line width=5pt] (190,60) -- (190,70);
\draw[color=black] (190,75) -- (190,113);
\draw[color=black, dash pattern={on 2.25pt off 2.25pt}] (190,113) -- (190,181);
\draw[color=black] (190,181) -- (190,208);

%BLACK points
\foreach \k in {-10,-8,-6,-4,-2,0,2,4,6,8,10} {
  \fill[black] ({190+12*\k}, {171-\k}) circle (2pt);
}

\foreach \k in {-7,7} {
  \fill[black] ({190+17*\k}, {171+\k}) circle (2pt);
}

% orange top polygon
\draw [color=orange, draw opacity=1, fill=orange, fill opacity=0.34]
  (235,93) -- (176,130) -- (158,115) -- cycle;

% gray polygon rest
\draw [draw opacity=0, fill=lightgray, fill opacity=0.17]
  (235,93) -- (207,185) -- (166,173) -- (158,115) -- cycle;


% Dashed gray lines
\draw[color=gray, dash pattern={on 2.25pt off 2.25pt}] (176,130) -- (204,154);
\draw[color=gray, dash pattern={on 2.25pt off 2.25pt}] (235,93) -- (204,154);

% gray left edges
\draw[color=gray] (158,115) -- (166,173);
\draw[color=gray] (176,130) -- (166,173);
\draw[color=gray] (176,130) -- (207,185);
\draw[color=gray] (235,93) -- (207,185);

% bottom darklav polygon
\draw [draw opacity=0, fill=darklav, fill opacity=0.11]
  (204,154) -- (207,185) -- (166,173) -- cycle;
\draw [color=darklav] (166,173) -- (207,185); %front edge
\draw [color=darklav, dash pattern={on 1.5pt off 1.5pt}] (204,154) -- (207,185); %dashed edges
\draw [color=darklav, dash pattern={on 1.5pt off 1.5pt}] (166,173) -- (204,154);


% DARKLAV VERTICES
\draw[color=darklav, fill=darklav] (190,171) circle (1.6pt); % center
\draw[color=darklav, fill=darklav] (207,185) circle (1.6pt); % front right
\draw[color=darklav, fill=darklav] (204,154) circle (1.6pt); % back right
\draw[color=darklav, fill=darklav] (166,173) circle (1.6pt); % front left



%ORANGE VERTICES
\draw[color=orange, fill=orange] (190,112) circle (1.6pt); % center 
\draw[color=orange, fill=orange] (158,115) circle (1.6pt); % top left
\draw[color=orange, fill=orange] (176,130) circle (1.6pt); % front
\draw[color=orange, fill=orange] (235,93)  circle (1.6pt); % back right

\end{scope}
\end{tikzpicture}


    \caption{The convex hull of two lattice triangles with each $\binom{2^2}{2}$ diameter directions which yields a $3$-dimensional lattice polytope with $\binom{2^3}{2}$ distinct diameter directions.}
    \label{fig:construction-diam-dir}
\end{figure}
%\input{Figures/3.1_diam-dir-triangles}
\end{example}


   



